\subsection{CLASSES OF EQUIVALENT OBJECTS}

Let \(\mathbb{R}\{x, y\}\) be an equivalence relation, which need not have a graph. It is obvious that if \(x, x', y\) are three distinct letters, the relation \(\mathbb{R}\{x, x'\}\) implies `` \(\mathbb{R}\{x, y\} \Longleftrightarrow \mathbb{R}\{x', y\}\) '' and therefore also implies the relation \((\forall y)(\mathbb{R}\{x, y\} \Longleftrightarrow \mathbb{R}\{x', y\})\) . By the scheme S7 (Chapter I, § 5, no. 1), if we put \(\theta\{x\} = \tau_y(\mathbb{R}\{x, y\})\) , the relation \(\mathbb{R}\{x, x'\}\) implies

\[
\theta \{x \} = \theta \left\{x ^ {\prime} \right\}.
\]

Note, on the other hand, that by definition \(\mathbb{R}\{x, \theta\{x\}\}\) is the relation \((\exists y)(\mathbb{R}\{x, y\})\) and hence (no. 1) is equivalent to \(\mathbb{R}\{x, x\}\) . It follows that the relation \((\mathbb{R}\{x, x\}\) and \(\mathbb{R}\{x', x'\}\) and \(\theta\{x\} = \theta\{x'\})\) is equivalent to \(\mathbb{R}\{x, x'\}\) , for it implies, by S6 (Chapter I, §5, no. 1), the relation

\[
(R \{x, x \} \text {   and   } R \{x ^ {\prime}, x ^ {\prime} \} \text {   and   } (R \{x ^ {\prime}, \theta \{x \} \} \iff R \{x ^ {\prime}, \theta \{x ^ {\prime} \} \})),
\]

hence also \((\mathbb{R}\{x, \theta\{x\}\}\) and \(\mathbb{R}\{x', \theta\{x\}\})\) , and finally \(\mathbb{R}\{x, x'\}\) by transitivity and symmetry. And since, conversely, \(\mathbb{R}\{x, x'\}\) implies \(\mathbb{R}\{x, x\}\) and \(\mathbb{R}\{x', x'\}\) , the assertion is proved. The term \(\theta\{x\}\) is called the class of objects equivalent to \(x\) (with respect to the relation \(\mathbb{R}\) ). Suppose now that \(\mathbb{T}\) is a term such that the relation

\[
(\forall y) (R \{y, y \} \Longrightarrow (\exists x) (x \in T \text {   and   } R \{x, y \}))\tag{1}
\]

is true. Then the relation \((\exists x)(R\{x, x\}\) and \(z = \theta\{x\})\) is collectivizing in \(z\) . For we may suppose that \(x \in T\) implies \(R\{x, x\}\) ; it is sufficient to replace \(T\) by the set of all \(x \in T\) such that \(R\{x, x\}\) (observing that \(R\{x, y\}\) implies \(R\{x, x\}\) ). Let \(\Theta\) be the set of all objects of the form \(\theta\{x\}\) for \(x \in T\) ( \(\S 1\) , no. 6). Suppose that \(R\{y, y\}\) is true; then there exists \(x \in T\) such that \(R\{x, y\}\) and therefore \(\theta\{y\} = \theta\{x\} \in \Theta\) . The set \(\Theta\) is called the set of classes of equivalent objects with respect to \(R\) ; and for each \(x\) such that \(R\{x, x\}\) , \(\theta\{x\}\) is the unique element \(z \in \Theta\) such that \(R\{x, z\}\) .

Under the same hypothesis, let \(\mathbf{A}\{x\}\) be a term such that \(\mathbf{R}\{x,y\}\) implies \(\mathbf{A}\{x\} = \mathbf{A}\{y\}\) . Then the relation \((\exists x)(\mathbf{R}\{x,x\}\) and \(z = \mathbf{A}\{x\})\) is collectivizing in \(z\) , since \(\mathbf{R}\{x,x\}\) , being equivalent to \(\mathbf{R}\{x,\theta\{x\}\}\) , implies \(\mathbf{A}\{x\} = \mathbf{A}\{\theta\{x\}\}\) , and consequently if \(\mathbf{E}\) is the set of objects of the form \(\mathbf{A}\{t\}\) for \(t \in \Theta\) , then \(\mathbf{R}\{x,x\}\) implies \(\mathbf{A}\{x\} \in \mathbf{E}\) . If \(f\) is the function \(t \to \mathbf{A}\{t\}\) ( \(t \in \Theta\) , \(\mathbf{A}\{t\} \in \mathbf{E}\) ), then the relation \(\mathbf{R}\{x,x\}\) implies \(\mathbf{A}\{x\} = f(\theta\{x\})\) .

In particular, if R is an equivalence relation on a set F, we may take A \(\{x\}\) to be the equivalence class of x with respect to R (no. 2), and the function f is then a bijection of \(\Theta\) onto the quotient set F/R; this justifies the terminology introduced.

* Example. Let \(\mathbb{R}\{x, y\}\) be the equivalence relation ``\(x\) and \(y\) are two vector spaces of the same finite dimension over \(\mathbf{C}\)''; this relation has no graph. It satisfies condition (1) when \(\mathbf{T}\) is the set of all vector subspaces of \(\mathbf{C}^{(\mathbf{N})}\) , or when the subset \(\mathbf{T}'\) of \(\mathbf{T}\) consists of the spaces \(\mathbf{C}^n\) ( \(n \in \mathbf{N}\) ) with the conventions that \(\mathbf{C}^0\) consists of the point 0 of \(\mathbf{C}^{(\mathbf{N})}\) and that \(\mathbf{C}^n\) ( \(n > 0\) ) is the sum of the first \(n\) components of the direct sum \(\mathbf{C}^{(\mathbf{N})}\) . With this second choice we have \(\Theta = \mathbf{T}'_*\)

